\section{Kernel moments for weighted finite-field compressions}
\label{sec:rank-moments}

The compact edition includes a finite-field estimate called
\path{card_mul_sum_det_le} in
\path{New/Heights/OneDeterminant.lean}.
It bounds singular matrix compressions using a positive-type weight on a
finite group, independently of prime distributions~\cite{repo}.
An earlier ProveIt report already extended its Fourier argument to ordered
frames and to the distribution of an entire rectangular
compression~\cite{priorSymplectic}. We retain that inherited argument and
extract a different observable: factorial moments of the kernel dimension.
These give direct rank-tail certificates which can remain useful when the
earlier total-variation certificate is vacuous.

The proofs below are complete human proofs. The additional statements have
not been compiled in Lean. Frame moments, or equivalently surjection moments
of the dual space, are established tools in random-matrix
theory~\cite{wood,fulman}; no worldwide priority claim is made for the
formulation here. The contribution of this section is an explicit transfer
of the repository's method, with its assumptions and quantitative costs
exposed.

\subsection{The positive-type estimate}

Let $V=\F_q^n$, let $G$ be a finite group, and let
$\pi:G\longrightarrow\operatorname{GL}(V)$ be a representation. The finite
field may have any characteristic. A probability mass function
$\mu:G\longrightarrow\mathbb R_{\ge0}$ is \emph{of positive type} if the
matrix
\[
 \bigl(\mu(g^{-1}h)\bigr)_{g,h\in G}
\]
is positive semidefinite over $\mathbb C$. This is an assumption in addition
to nonnegativity and total mass one. Write $g\cdot x=\pi(g)x$ and
$g\cdot y=y\circ\pi(g)^{-1}$ for the actions on $V$ and $V^*$. Define
\[
 \alpha(x)=\mu\{g:g\cdot x=x\},\qquad
 \beta(y)=\mu\{g:g\cdot y=y\}.
\]
Fix a nontrivial additive character
$\psi:\F_q\longrightarrow\mathbb C^\times$.

\begin{lemma}[Positive-type coefficient bound]
\label{lem:positive-coefficient}
For every $x\in V$ and $y\in V^*$,
\[
 \left|\E_\mu\psi\bigl(y(\pi(g)x)\bigr)\right|
 \le q^{n/2}\sqrt{\alpha(x)\beta(y)}.
\]
\end{lemma}

\begin{proof}
Choose vectors $u_g$ in a complex Hilbert space with
$\langle u_g,u_h\rangle=\mu(g^{-1}h)$; a finite positive semidefinite
matrix always has such a Gram representation. Inner products are
conjugate-linear in the first argument. In $\ell^2(V)$ let $e_x$ be the
point masses and put
\[
 c_y(z)=q^{-n/2}\psi(-y(z)).
\]
Additive-character orthogonality makes both indexed families orthonormal.
Consider the two finite tensor sums
\[
 P=\sum_{g\in G}c_{g\cdot y}\otimes u_g,
 \qquad
 Q=\sum_{h\in G}e_{h\cdot x}\otimes u_h.
\]
Setting $t=g^{-1}h$ in their inner product gives
\[
 \langle P,Q\rangle
 =|G|q^{-n/2}\sum_{t\in G}\mu(t)\psi\bigl(y(\pi(t)x)\bigr).
\]
Orthonormality similarly gives
$\|P\|^2=|G|\beta(y)$ and $\|Q\|^2=|G|\alpha(x)$.
Cauchy--Schwarz proves the result.
\end{proof}

Applying the lemma on $V^k$ gives the version for frames. If
$X=(x_1,\ldots,x_k)$ and $Y=(y_1,\ldots,y_k)$, then
\begin{equation}
 \left|\E_\mu\psi\!\left(\sum_{i=1}^k y_i(\pi(g)x_i)\right)\right|
 \le q^{nk/2}
       \sqrt{\mu(\operatorname{Stab}X)\,
                    \mu(\operatorname{Stab}Y)}.
 \label{eq:frame-coefficient}
\end{equation}
Here the stabilizers fix every member of the frame; the second uses the
dual action. This extension itself is inherited from the earlier
report~\cite{priorSymplectic}. No transitivity or irreducibility is needed.

\subsection{An exact first moment and its rank tail}

Let $A:\F_q^a\to V$ and $B:\F_q^b\to V^*$ be linear maps, initially
not required to be injective. Define the $b\times a$ matrix $M(g)$ by
\[
 M(g)_{ij}=(Be_i)(\pi(g)Ae_j),\qquad
 \nu(g)=\dim_{\F_q}\ker M(g).
\]
In coordinate identifications this is $B^{\mathsf T}\pi(g)A$. Put
\[
 P_A=\sum_{v\in\F_q^a\setminus\{0\}}\sqrt{\alpha(Av)},
 \qquad
 P_B=\sum_{y\in\F_q^b\setminus\{0\}}\sqrt{\beta(By)}.
\]

\begin{proposition}[Kernel-size moment]
\label{prop:kernel-first-moment}
The exact independent-entry baseline and an error certificate are
\begin{equation}
 \left|\E_\mu q^{\nu(g)}-
      \left(1+\frac{q^a-1}{q^b}\right)\right|
 \le q^{n/2-b}P_AP_B.
 \label{eq:first-kernel-moment}
\end{equation}
Consequently, for $1\le r\le a$,
\begin{equation}
 \Prob_\mu(\nu\ge r)
 \le \min\!\left\{1,
 \frac{(q^a-1)q^{-b}+q^{n/2-b}P_AP_B}{q^r-1}\right\}.
 \label{eq:first-kernel-tail}
\end{equation}
\end{proposition}

\begin{proof}
Count the kernel by additive-character orthogonality:
\[
 \E_\mu q^\nu
 =q^{-b}\sum_{v\in\F_q^a}\sum_{y\in\F_q^b}
       \E_\mu\psi\bigl((By)(\pi(g)Av)\bigr).
\]
Pairs for which $v=0$ or $y=0$ contribute exactly
$q^{-b}(q^a+q^b-1)$. Bound each remaining coefficient with
Lemma~\ref{lem:positive-coefficient}, and factor the two finite sums of square
roots. This proves the two-sided moment estimate. The nonnegative variable
$q^\nu-1$ is at least $q^r-1$ on $\{\nu\ge r\}$; Markov's inequality gives
the tail estimate.
\end{proof}

For square compressions $a=b=m$, if every relevant nonzero-vector
stabilizer mass is at most $Y$, the moment estimate becomes
\[
 \left|\E_\mu q^\nu-(2-q^{-m})\right|
 \le q^{m+n/2}(1-q^{-m})^2Y.
\]
The repository's displayed singular-minor conclusion instead has the form
$q\Prob(\det M=0)\le2+q^{m+n/2}Y$. Keeping the zero-character contribution
exact and retaining the kernel-size moment gives information about every
deficiency threshold, rather than just singularity. The final denominator
$q^r-1$ is best possible if only the expectation of $q^\nu-1$ is known:
an abstract variable supported on $\{0,r\}$ attains equality in that Markov
step. This does not claim simultaneous sharpness of the preceding Fourier
estimate.

\subsection{A hierarchy of factorial kernel moments}

Assume henceforth that $A$ and $B$ are injective, with $a,b\le n$.
For nonnegative integers $d,j$, let
\[
 F_q(d,j)=
 \begin{cases}
  \displaystyle\prod_{i=0}^{j-1}(q^d-q^i),&0\le j\le d,\\[3pt]
  0,&j>d.
 \end{cases}
\]
Thus $F_q(d,j)$ counts ordered independent $j$-frames in $\F_q^d$,
and $F_q(d,0)=1$. The number of $b\times j$ matrices of rank $k$ is
\begin{equation}
 R_q(b,j;k)=\frac{F_q(b,k)F_q(j,k)}{F_q(k,k)}.
 \label{eq:rank-matrix-count}
\end{equation}
Indeed, factor a rank-$k$ matrix as a full-column-rank $b\times k$ matrix
times a full-row-rank $k\times j$ matrix, and divide by the
$|\operatorname{GL}_k(\F_q)|$ changes of intermediate basis.

Let $\alpha_k$ be the maximum $\mu$-mass of the pointwise stabilizer of
an independent $k$-frame in $\operatorname{im}A$. Let $\beta_k$ be the
corresponding maximum for independent $k$-frames in $\operatorname{im}B$
under the dual action. For $1\le j\le a$, define
\begin{align}
 U_j&=\frac{F_q(a,j)}{q^{bj}},\label{eq:moment-main}\\
 \varepsilon_j
   &=\sum_{k=1}^{\min(b,j)}R_q(b,j;k)\,
          q^{nk/2}\sqrt{\alpha_k\beta_k}.
          \label{eq:moment-epsilon}
\end{align}

\begin{theorem}[Factorial kernel moments]
\label{thm:factorial-kernel}
For every $1\le j\le a$,
\begin{equation}
 \left|\E_\mu F_q(\nu(g),j)-U_j\right|
 \le U_j\varepsilon_j.
 \label{eq:factorial-moment}
\end{equation}
For every $1\le j\le r\le a$ it follows that
\begin{equation}
 \Prob_\mu(\nu\ge r)
 \le \min\!\left\{1,
       \frac{U_j(1+\varepsilon_j)}{F_q(r,j)}\right\}.
 \label{eq:factorial-tail}
\end{equation}
In particular, one may minimize the right-hand side over $j$.
\end{theorem}

\begin{proof}
For each $g$, the quantity $F_q(\nu(g),j)$ counts full-column-rank matrices
$T\in\operatorname{Mat}_{a\times j}(\F_q)$ for which $M(g)T=0$.
For a fixed such $T$, character orthogonality on all $b\times j$ matrices
gives
\[
 \boldsymbol{1}_{M(g)T=0}
 =q^{-bj}\sum_{W\in\operatorname{Mat}_{b\times j}(\F_q)}
       \psi\bigl(\operatorname{tr}(W^{\mathsf T}M(g)T)\bigr).
\]
Take expectations and sum over the $F_q(a,j)$ choices of $T$.
The term $W=0$ contributes exactly $U_j$.

For a nonzero $W$ of rank $k$, choose a factorization $W=CD$, where $C$
has size $b\times k$ and rank $k$, and $D$ has size $k\times j$ and
rank $k$. Cyclicity of the trace yields
\[
 \operatorname{tr}(W^{\mathsf T}M(g)T)
 =\operatorname{tr}(C^{\mathsf T}M(g)TD^{\mathsf T}).
\]
The columns of $ATD^{\mathsf T}$ form an independent $k$-frame in
$\operatorname{im}A$: each of $D^{\mathsf T}$, $T$, and $A$ is injective.
The columns of $BC$ form an independent $k$-frame in $\operatorname{im}B$.
Equation~\eqref{eq:frame-coefficient} therefore bounds the absolute value
of the expected character by $q^{nk/2}\sqrt{\alpha_k\beta_k}$.
This bound is independent of $T$ and of the chosen factorization.

There are $R_q(b,j;k)$ possibilities for $W$ of rank $k$.
The triangle inequality, summed over $W$ and $T$, proves
\eqref{eq:factorial-moment}. Finally, $F_q(\nu,j)$ is nonnegative and is
at least $F_q(r,j)$ when $\nu\ge r$. Apply Markov's inequality and the
upper moment bound to obtain~\eqref{eq:factorial-tail}.
\end{proof}

The normalization has an exact probabilistic meaning. If $Z$ is uniform
on all $b\times a$ matrices, annihilating any fixed independent $j$-frame
imposes $bj$ independent linear equations. Hence
\[
 \E F_q(\dim\ker Z,j)=F_q(a,j)q^{-bj}=U_j.
\]
These are precisely the factorial kernel moments of the independent-entry
model, rather than an analogy based only on the mean rank. Counting frames
is equivalent, by duality, to counting surjections to $\F_q^j$; that moment
viewpoint has well-developed precedents~\cite{wood}. Exact finite-field rank
laws and total-variation comparisons also have an established
literature~\cite{fulman}.

\subsection{Comparison with the earlier full-distribution bound}

The prior report obtains the certificate
\begin{equation}
 d_{\mathrm{TV}}(M(g),Z)
 \le\frac12\left(
   \sum_{k=1}^{\min(a,b)}R_q(b,a;k)q^{nk}\alpha_k\beta_k
       \right)^{1/2}.
 \label{eq:inherited-TV}
\end{equation}
It controls all matrix events simultaneously~\cite{priorSymplectic}.
The new moment bound instead sums over frequency matrices of size
$b\times j$, not $b\times a$. It requires frame stabilizers only through
rank $\min(b,j)$ and can focus on a particular rank tail. The two bounds
use different norms and quantities; neither dominates the other in all
regimes. A small total-variation bound can be far stronger. The following
example shows why the direct certificate is nevertheless useful.

Take $G=\operatorname{GL}_{32}(\F_2)$ with uniform probability, $a=31$,
$b=4$, and coordinate injections $A,B$. The event
$\nu\ge28$ is exactly failure of full row rank. Transitivity on independent
frames, also for the dual action, gives
\[
 \alpha_k=\beta_k=\frac{1}{F_2(32,k)}.
\]
For $j=8$, exact rational arithmetic gives
$\varepsilon_8<0.058894$. The resulting certificates are as follows.
All displayed decimal upper bounds are rounded upward.

\begin{center}
\begin{tabular}{p{0.55\linewidth}p{0.33\linewidth}}
\hline
Quantity & Value or certificate\\
\hline
First moment, $j=1$ & $\Prob(\nu\ge28)\le0.500115$\\
Eighth moment, $j=8$ & $\Prob(\nu\ge28)\le0.004137$\\
Uncapped expression in~\eqref{eq:inherited-TV}
 & $>2.132$, hence vacuous\\
Exact probability in this particular $GL$ example
 & $15/(2^{32}-1)$\\
\hline
\end{tabular}
\end{center}

The exact eighth-moment upper bound, before decimal rounding, is
\[
 \frac{1051362055003754754975737756578755}
      {254178960140462084616340665075236864}.
\]
The true probability is below $3.492460\cdot10^{-9}$, so this example
does \emph{not} improve the specialized exact $GL_n$ calculation.
To see the exact value, $\pi(g)\operatorname{im}A$ is a uniform hyperplane
in $\F_2^{32}$. Projection to $\F_2^4$ loses one row of rank exactly when
that hyperplane contains the fixed $28$-dimensional projection kernel.
There are $2^{32}-1$ hyperplanes in total and $2^4-1$ containing this kernel.
No larger loss is possible. The useful comparison is between general
certificates using the same stabilizer information. In actions or weighted
laws without an explicit compression distribution, the direct rank bound
can supply information when~\eqref{eq:inherited-TV} supplies none.

\subsection{Assumptions and exact finite checks}

Positive type is essential, even for a nonnegative, normalized,
inverse-symmetric weight. Take $G=\operatorname{GL}_2(\F_2)$, $a=b=1$,
$A=B=e_1$, and put all probability on
\[
 s=\begin{pmatrix}0&1\\1&0\end{pmatrix}.
\]
Both relevant stabilizer masses are zero, but $M(s)=0$ with probability
one. Thus $\E 2^\nu=2$, whereas the center in
\eqref{eq:first-kernel-moment} is $3/2$ and its asserted error would be
zero. This law is not of positive type: the principal Gram matrix indexed
by $\{1,s\}$ is $\left(\begin{smallmatrix}0&1\\1&0\end{smallmatrix}\right)$,
with determinant $-1$. Conversely, positive type alone does not make a
general complex-valued function a probability; nonnegativity is also needed
for the expectation and Markov arguments.

The supporting Python script enumerates all $20{,}160$ elements of
$\operatorname{GL}_4(\F_2)$ and the $720$ elements of its standard symplectic
subgroup. It checks square $2\times2$ and rectangular $2\times3$
compressions for two laws: uniform $GL_4$, and the nonuniform mixture
\[
 \tfrac12\operatorname{Unif}(GL_4)
 +\tfrac12\operatorname{Unif}(Sp_4).
\]
Both laws are of positive type. For any subgroup $H$, the kernel
$\boldsymbol{1}_H(g^{-1}h)/|H|$ is a positive semidefinite block matrix on
right cosets; convex combinations preserve positivity. In the mixture the
exact weights are $29/40320$ on $Sp_4$ and $1/40320$ elsewhere.

All $636$ relevant ordered-frame Fourier inequalities are checked after
squaring, avoiding floating-point square roots. The script independently
counts annihilated frames, verifies the independent-entry main terms, checks
every stated finite moment and tail, and uses exact rational arithmetic for
the large $GL$ examples. For the square compression the two computed
nullity laws are
\[
\begin{array}{c|ccc}
 &\nu=0&\nu=1&\nu=2\\ \hline
 \operatorname{Unif}(GL_4)&16/35&18/35&1/35\\
 \tfrac12\operatorname{Unif}(GL_4)+
 \tfrac12\operatorname{Unif}(Sp_4)&52/105&16/35&1/21
\end{array}
\]
There are $1{,}277$ successful checks, including bookkeeping and exact
rounding checks. The JSON output records the exact fractions and the
verification scope. These experiments verify finite examples of the
theorems proved above; they do not replace the proofs, establish novelty,
or constitute a Lean replay of the repository.
